\section{Frequency ranges and positivity}
\label{sec:frequency-positivity}

After fibre compression, it remains to estimate the decoded frequency $h_m$ attached to each coherent anchor label $m$. Noncoherent anchor assignments and nondecoder row tuples already have the bounds \cref{eq:noncoherent-top-tail,eq:global-fibre-error}. The retained prime-pair factors provide the remaining decay.

\subsection{Elementary summation bounds}

\begin{lemma}[Gaussian and stretched-exponential tails]
\label{lem:tail-summation}
Fix $c>0$ and a nonnegative integer $d$.
\begin{enumerate}[label=(\roman*)]
\item There are constants $C,c'>0$ such that, for $B\geq1$ and $A\geq B$,
\[
 \sum_{m\geq A}(1+m/B)^d\exp(-cm^2/B^2)
 \leq CB(1+A/B)^d\exp(-c'A^2/B^2).
\]
\item There are constants $C,c'>0$ and $A_0$ such that, for $A\geq A_0$,
\[
 \sum_{m\geq A}m^d\exp\left(-\frac{cm}{(\log m)^2}\right)
 \leq C\exp\left(-\frac{c'A}{(\log A)^2}\right).
\]
\end{enumerate}
\end{lemma}
\begin{proof}
For~(i), compare with the corresponding integral, put $u=x/B$, and use the standard Gaussian tail bound. For~(ii), split into dyadic intervals. On the $j$th interval the exponential decay absorbs the polynomial length and weight, and the resulting geometric series is dominated by its first term.
\end{proof}

\subsection{The prime-only transition}

For $T_0<|m|\leq X^2/4$, every coordinate indexed by $P$ equals $m$. Hence every complete-pair factor has phase $m/(pq)$ in the linear circle range, and
\begin{equation}
\label{eq:transition-pointwise}
 |F(h_m)|\leq\exp\left(-c\frac{m^2}{X^2\log^2X}\right).
\end{equation}
Moreover,
\[
 \frac{T_0^2}{X^2\log^2X}\longrightarrow\infty
\]
for every fixed $\gamma>1$. Applying \cref{lem:tail-summation}(i),
\begin{equation}
\label{eq:transition-total}
 \sum_{T_0<|m|\leq X^2/4}|F(h_m)|=o(1)=o(\sigma_{\mathcal E}^{-1}).
\end{equation}
This range uses no target coordinate.

\subsection{Adaptive retained-pair damping}

For $X^2/4<|m|\leq M_{\rm dec}$ put
\[
 I_m=[2\sqrt{|m|},3\sqrt{|m|}].
\]

\begin{proposition}[Adaptive prime interval]
\label{prop:adaptive-prime-interval}
For every such $m$, the interval $I_m$ lies in $[X,\eta Z)$ for all sufficiently large $X$ and contains
\[
 K_m\gg\frac{\sqrt{|m|}}{\log|m|}
\]
primes. The retained prime-pair energy satisfies
\begin{equation}
\label{eq:adaptive-energy}
 Q_{\rm pair}(m)\gg\frac{|m|}{(\log|m|)^2}.
\end{equation}
\end{proposition}
\begin{proof}
The lower inclusion follows from $2\sqrt{|m|}>X$. At the upper endpoint,
\[
 3\sqrt{|m|}\leq\frac{3\sqrt{XZ}}{\log Z}=o_\eta(Z)<\eta Z.
\]
The prime number theorem in a fixed-ratio interval gives the count. For distinct primes $p,q\in I_m$, one has $1/9\leq|m|/(pq)\leq1/4$, so every retained pair contributes a fixed positive amount of circle energy. Summing over the pairs proves \cref{eq:adaptive-energy}.
\end{proof}

Thus
\[
 |F(h_m)|\leq\exp\left(-c\frac{|m|}{(\log|m|)^2}\right),
\]
and \cref{lem:tail-summation}(ii) gives
\begin{equation}
\label{eq:adaptive-total}
 \sum_{X^2/4<|m|\leq M_{\rm dec}}|F(h_m)|
 \ll\exp\left(-c'\frac{X^2}{\log^2X}\right)
 =o(1)=o(\sigma_{\mathcal E}^{-1}).
\end{equation}

\subsection{The six exhaustive ranges}

Fix $C\geq1$, to be chosen before $X$, and put
\begin{equation}
\label{eq:N-major}
 N=\left\lfloor\frac{C}{\sigma_{\mathcal E}}\right\rfloor.
\end{equation}
For fixed $C$, the total variance gives
\begin{equation}
\label{eq:N-below-full}
 N=o(T_0).
\end{equation}
Thus, for sufficiently large $X$,
\[
 N<T_0<\frac{X^2}{4}<M_{\rm dec}.
\]
The decoder/nondecoder split is made first. Every nondecoder row tuple belongs to Range~VI. Among decoded tuples, every noncoherent anchor assignment also belongs to Range~VI. A coherent assignment has one unique integer label $m$ and enters exactly one of the following disjoint ranges:
\begin{center}
\begin{tabular}{@{}c p{5.2cm} p{5.0cm}@{}}
\toprule
Range & condition & arithmetic control \\
\midrule
I & $|m|\leq N$ & positive Taylor main term \\
II & $N<|m|\leq T_0$ & total-variance Gaussian decay \\
III & $T_0<|m|\leq X^2/4$ & prime-pair Gaussian decay \\
IV & $X^2/4<|m|\leq M_{\rm dec}$ & adaptive retained pairs \\
V & $|m|>M_{\rm dec}$ & internal anchor energy \\
VI & noncoherent decoded assignment, or a nondecoder row tuple & anchor counting and fibre error \\
\bottomrule
\end{tabular}
\end{center}
These cases are pairwise disjoint and exhaustive because the CRT coordinate map is a bijection, each row fibre has one chosen decoder tuple, and every coherent anchor assignment has exactly one label.

For Range~V, the anchor factor and \cref{eq:top-exact-energy} give
\[
 |F(h_m)|\leq\exp(-\kappa_{\rm mod}m^2\sigma_{B,0}^2).
\]
Since $M_{\rm dec}\sigma_{B,0}\asymp_\eta X/\log^3Z\to\infty$, \cref{lem:tail-summation}(i) yields
\begin{equation}
\label{eq:anchor-tail-total}
 \sum_{|m|>M_{\rm dec}}|F(h_m)|
 \ll_\eta Z\log Z\exp\left(-c_\eta\frac{X^2}{\log^6Z}\right)
 =o(1)=o(\sigma_{\mathcal E}^{-1}).
\end{equation}
Range~VI has total modulus
\begin{equation}
\label{eq:terminal-error-total}
 \ll_\eta\exp(-c_\eta F_B)+P_{\rm anc}(\exp(\Delta)-1)
 =o(1)=o(\sigma_{\mathcal E}^{-1}).
\end{equation}

\subsection{The central range}

Compactness of the Bernoulli interval gives constants $\rho>0$ and $M<\infty$ such that, uniformly for $\theta\in I_{t,\gamma,\eta}$ and $|z|\leq\rho$,
\begin{equation}
\label{eq:local-log-expansion}
 \log\bigl((1-\theta)+\theta\exp(2\pi iz)\bigr)
 =2\pi i\theta z-2\pi^2\theta(1-\theta)z^2+R_\theta(z),
 \qquad |R_\theta(z)|\leq M|z|^3.
\end{equation}
Let $e_{\min}=\min_{e\in\mathcal E}e\asymp_{b,\eta}\min(X^2,Z)$. For $|m|\leq N$,
\[
 \sum_{e\in\mathcal E}|m/e|^3
 \leq\frac{|m|^3}{e_{\min}}\sum_{e\in\mathcal E}\frac1{e^2}
 \ll_{b,\eta}\frac{C^3}{\sigma_{\mathcal E}e_{\min}}=o(1).
\]
Also $N=o(e_{\min})$, so every phase lies in the Taylor disk.

For $|m|\leq N$, \cref{eq:h-m-full-identification} identifies the decoded frequency with the integer $m$. The linear phase cancels exactly by \cref{eq:exact-centering-load}; the quadratic term is the total actual variance. Uniformly in this range,
\begin{equation}
\label{eq:major-local-gaussian}
 F(m)=\exp\bigl(-2\pi^2m^2\sigma_{\mathcal E}^2+\varepsilon_m\bigr),
 \qquad \max_{|m|\leq N}|\varepsilon_m|\to0.
\end{equation}
It follows that
\begin{proposition}[Positive Gaussian contribution]
\label{prop:positive-major}
For every fixed $C\geq1$ and all sufficiently large $X$,
\begin{equation}
\label{eq:positive-major}
 \RePart\sum_{|m|\leq N}F(m)\geq\frac{c_{\rm maj}}{\sigma_{\mathcal E}},
\end{equation}
where $c_{\rm maj}>0$ is independent of $C$ and $X$.
\end{proposition}

\subsection{The first Gaussian tail and final assembly}

For $N<|m|\leq T_0$, every coordinate is again the integer $m$, and $|m|/e\leq1/4$ for all $e\in\mathcal E$. Hence
\[
 |F(m)|\leq\exp(-cm^2\sigma_{\mathcal E}^2).
\]
The Gaussian tail estimate gives
\begin{equation}
\label{eq:range-II-tail}
 \sum_{|m|>N}\exp(-cm^2\sigma_{\mathcal E}^2)
 \leq K\sigma_{\mathcal E}^{-1}\exp(-c'C^2).
\end{equation}
Choose $C$ so large that this is at most $c_{\rm maj}/(8\sigma_{\mathcal E})$, and then take $X$ large enough for all preceding estimates.

The complete Fourier sum is the sum of the positive Range~I contribution, the decoded contributions in Ranges~II--V, and the Range~VI error. By the triangle inequality, \cref{eq:transition-total,eq:adaptive-total,eq:anchor-tail-total,eq:terminal-error-total,eq:range-II-tail} and \cref{eq:positive-major} give
\[
 \RePart\sum_{h\bmod L}F(h)
 \geq\frac{c_{\rm maj}}{\sigma_{\mathcal E}}
 -\frac{c_{\rm maj}}{8\sigma_{\mathcal E}}
 -o(\sigma_{\mathcal E}^{-1})>0.
\]
We have proved:

\begin{theorem}[Positive Fourier numerator]
\label{thm:positive-fourier-numerator}
For every fixed reduced $t=a/b\in(0,1)$ with squarefree $b$, every fixed $\eta\in(0,1)$, every fixed $\gamma$ satisfying \cref{eq:admissible-region}, and every finite forbidden set $\mathcal T$, one has, for all sufficiently large $X$,
\[
 \RePart\sum_{h\bmod L}F(h)>0.
\]
\end{theorem}
